%------------------------------------------------------------------------------%
\documentclass[fleqn,utf8,aspectratio=169,12pt]{beamer}

\usepackage{gaal}

\title{Produto por Escalar}

%------------------------------------------------------------------------------%
\begin{document}

\addtitle

%------------------------------------------------------------------------------%
\section{Produto por Escalar}

%------------------------------------------------------------------------------%
\begin{frame}{Produto de um Vetor por um Escalar}
  O produto de $\vec{u} \in \R^n$
  \pause
  por $\lambda \in \R$ é um novo vetor
  \[
    \pause
    \vec{v} = \lambda \vec{u}
  \]

  \pause
  O vetor $\vec{u}$ é escalado pelo valor $\lambda$

  \pause
  O vetor $\vec{v}$ é um múltiplo do vetor $\vec{u}$

  \pause
  O vetor $\vec{v}$ é colinear, ou paralelo, ao vetor $\vec{u}$
\end{frame}

%------------------------------------------------------------------------------%
\framedgraphic*{Interpretação Geométrica}{E-02-Produto_por_escalar-1}{}
\framedgraphic*{Interpretação Geométrica}{E-02-Produto_por_escalar-2}{}
\framedgraphic*{Interpretação Geométrica}{E-02-Produto_por_escalar-3}{}

%------------------------------------------------------------------------------%
\begin{frame}{Produto de um Vetor por um Escalar}
  \pause
  Se \; $\lambda>1$   \tabto{30mm} o vetor aumenta de comprimento

  \pause
  \smallskip
  Se \; $0<\lambda<1$ \tabto{30mm} o vetor diminui de comprimento

  \pause
  \smallskip
  Se \; $\lambda<0$   \tabto{30mm} o sentido é invertido

  \pause
  \smallskip
  Se \; $\lambda=0$   \tabto{30mm} obtemos o vetor nulo
\end{frame}

%------------------------------------------------------------------------------%
\begin{frame}{Cálculo do Produto por Escalar}
  Para
  \(
    \lambda \in \R
  \)
  \pause
  e
  \(
    \vec{u}
    = ( u_1, u_2 )
    \in \R^2
  \)

  \pause
  \[
    \lambda\vec{u}
    \pause
    = \lambda \left( u_1, u_2 \right)
    \pause
    = \left( \lambda u_1, \lambda u_2 \right)
    \pause
    = \vetor{\lambda u_1}{\lambda u_2}
    \pause
    = \lambda \vetor{u_1}{u_2}
  \]
\end{frame}

%------------------------------------------------------------------------------%
\input{ex/E-02-Produto_por_escalar-ex-1}

%------------------------------------------------------------------------------%
\begin{frame}{Propriedades do Produto por Escalar}
  Para escalares $\lambda$ e $\mu$ e vetores $\vec{u}$ e $\vec{v}$
  \[
    \pause
    1 \, \vec{u} = \vec{u}
  \]

  \[
    \pause
    0 \, \vec{u} = \vec{0}
  \]

  \[
    \pause
    (-1)\vec{u} = -\vec{u}
  \]

  \[
    \pause
    \lambda(\mu\vec{u}) = (\lambda\mu)\vec{u}
  \]

  \[
    \pause
    (\lambda+\mu)\vec{u} = \lambda\vec{u}+\mu\vec{u}
  \]

  \[
    \pause
    \lambda(\vec{u}+\vec{v}) = \lambda\vec{u}+\lambda\vec{v}
  \]
\end{frame}

%------------------------------------------------------------------------------%
%------------------------------------------------------------------------------%
\begin{question}
  Determine $\lambda$ sabendo que
  \(
    \lambda \vetor{2}{-3} = \vetor{8}{-12}
  \)
\end{question}


%------------------------------------------------------------------------------%
\begin{resolution}{Solução}
  \begin{align*}
    \onslide<+->{\lambda \vetor{2}{-3}         & = \vetor{8}{-12} \\}
    \onslide<+->{\vetor{2\lambda }{-3\lambda } & = \vetor{8}{-12}}
  \end{align*}

  \onslide<+->{\[
    \left\{\begin{array}{r<{\!\!\!}l}
      2\lambda  & = 8 \\
      -3\lambda & = -12
    \end{array}
    \right.
  \]}
\end{resolution}

%------------------------------------------------------------------------------%
\begin{resolution}{Solução}
  \onslide<+->{Comparando a primeira componente}
  \begin{align*}
    \onslide<+->{2\lambda & = 8 \\}
    \onslide<+->{\lambda  & = 4}
  \end{align*}
  \onslide<+->{Verificando a segunda componente}
  \begin{align*}
    \onslide<+->{-3\lambda & = -12 \\}
    \onslide<+->{\lambda   & = 4}
  \end{align*}
  \onslide<+->{Portanto
  \[
    \lambda = 4
  \]}
\end{resolution}

%------------------------------------------------------------------------------%
\endinput



%------------------------------------------------------------------------------%
\begin{question}
  Determine $p$ para que os vetores \(\vec{v}\) e \(\vec{w}\) sejam paralelos
  \[
    \vec{v} = (-6, 3, p)
  \]
  \[
    \vec{w} = ( 2, -1, 3 )
  \]
\end{question}

%------------------------------------------------------------------------------%
\begin{resolution}{Solução}
  \onslide<+->{Os vetores serão paralelos se um for múltiplo do outro,}
  \onslide<+->{isto é, se \emph{existir $\lambda$} tal que}
  \begin{align*}
    \onslide<+->{\vec{v}    & = \lambda\vec{w}       \\}
    \onslide<+->{(-6, 3, p) & = \lambda ( 2, -1, 3 ) \\}
    \onslide<+->{(-6, 3, p) & = \big( \lambda 2,\, \lambda (-1),\, \lambda 3 \big)\\}
    \onslide<+->{(-6, 3, p) & = \big( 2\lambda,\, -\lambda,\, 3\lambda \big)}
  \end{align*}
\end{resolution}

%------------------------------------------------------------------------------%
\begin{resolution}{Solução}
  \begin{columns}[t]
  \column{0.25\linewidth}
    \begin{align*}
      \onslide<+->{-6       & = 2\lambda \\}
      \onslide<+->{2\lambda & = -6       \\}
      \onslide<+->{\lambda  & = -3}
    \end{align*}
  \column{0.4\linewidth}
    \begin{align*}
      \onslide<+->{3        & = -\lambda \\}
      \onslide<+->{-\lambda & = 3        \\}
      \onslide<+->{\lambda  & = -3}
    \end{align*}
    \qquad\onslide<+->{compatível}
  \column{0.35\linewidth}
    \begin{align*}
      \onslide<+->{p & = 3\lambda \\}
      \onslide<+->{  & = 3(-3)    \\}
      \onslide<+->{  & = -9}
    \end{align*}
  \end{columns}
\end{resolution}

%------------------------------------------------------------------------------%
\endinput


%------------------------------------------------------------------------------%
%\section{Lista Mínima}
%
%%------------------------------------------------------------------------------%
%\begin{frame}{Lista Mínima}
%  \begin{enumerate}
%    \item
%  \end{enumerate}
%
%  \vfill
%  Atenção: A prova é baseada no livro, não nas apresentações
%\end{frame}

%------------------------------------------------------------------------------%
\end{document}
%------------------------------------------------------------------------------%


